\documentclass{article}
\usepackage{amsmath,amssymb,xurl}
\usepackage[hidelinks]{hyperref}
% Shared commands for notes in docs/paper-gaps/.

\DeclareUnicodeCharacter{2080}{\ifmmode{}_{0}\else\textsubscript{0}\fi}
\DeclareUnicodeCharacter{2081}{\ifmmode{}_{1}\else\textsubscript{1}\fi}
\DeclareUnicodeCharacter{2082}{\ifmmode{}_{2}\else\textsubscript{2}\fi}
\DeclareUnicodeCharacter{2099}{\ifmmode{}_{n}\else\textsubscript{n}\fi}

\newcommand{\Lean}{\textsc{Lean}}
\ExplSyntaxOn
\NewDocumentCommand{\leanid}{m}{
  \ifmmode
    \textsf{\detokenize{#1}}
  \else
    \group_begin:
    \urlstyle{sf}
    \tl_set:Nn \l_tmpa_tl {#1}
    \tl_replace_all:Nnn \l_tmpa_tl {\_} {_}
    \exp_args:NV \path \l_tmpa_tl
    \group_end:
  \fi
}
\ExplSyntaxOff
\providecommand{\tr}{\operatorname{tr}}
\newcommand{\tnleanIssueBase}{https://github.com/LionSR/TNLean/issues/}
\newcommand{\tnleanPrBase}{https://github.com/LionSR/TNLean/pull/}
\newcommand{\tnleanDocBase}{https://sirui-lu.com/TNLean/docs/}
\newcommand{\ghissue}[1]{\href{\tnleanIssueBase#1}{GitHub issue \##1}}
\newcommand{\ghpr}[1]{\href{\tnleanPrBase#1}{PR \##1}}
\newcommand{\doclink}[1]{\href{\tnleanDocBase#1.html}{\path{#1}}}

% Dirac notation and the identity operator, as in blueprint/src/macros/common.tex.
% \providecommand keeps note-local definitions authoritative.
\usepackage{dsfont}
\providecommand{\ket}[1]{|#1\rangle}
\providecommand{\bra}[1]{\langle #1|}
\providecommand{\braket}[2]{\langle #1 | #2 \rangle}
\providecommand{\ketbra}[2]{|#1\rangle\!\langle #2|}
\providecommand{\Id}{\mathds{1}}
\providecommand{\one}{\mathds{1}}
\providecommand{\C}{\mathbb{C}}
\providecommand{\MN}[1]{M_{#1}(\C)}
\providecommand{\MD}{M_D(\C)}

% Verdict marker: \gapnote{<kind>}{<status>}. Kind is the policy
% classification (clarification, local-correction, scope-restriction,
% unfaithful, false-source, open-gap); status is open, wip (elimination
% underway), resolved, or historical. Machine-read by the site index;
% prints a verdict line.
\newcommand{\gapnote}[2]{\par\noindent\textbf{Verdict:} #1 (#2).\par}

\title{Native Open Central Flux Pairs and Vacuum Reduced Density}
\author{}
\date{}
\begin{document}
\maketitle
\gapnote{scope-restriction}{open}

Schuch, Cirac, and P\'erez-Garc\'ia,
\href{https://arxiv.org/abs/1001.3807v3}{arXiv:1001.3807v3},
Definition 6.13 and Lemma 6.14, source
\path{Papers/1001.3807/paper_v3.tex}, lines 2180--2214,
construct flux pairs and assert local unobservability of the connecting string.

For a physical observation set $R$, let $B(R)$ be the explicitly computed
set of dual plaquettes with at least one corner in $R$. Assume both periods
are at least three, the actual closed-cell realization of $B(R)$ is simply
connected, and both endpoints of a finite native dual path $p$ lie outside
$B(R)$. Let the original tensor be regular $G$-isometric and the insertion
element $g$ be central. The theorem
\leanid{TNLean.PEPS.IsGIsometric.torusCentralFlux_normalized_density_eq_vacuum}
proves that the actual flux and vacuum cut matrices are nonzero and have
identical unnormalized reduced matrices, hence identical normalized reduced
densities. The endpoints need not coincide and the winding is unrestricted.

The proof derives an exterior detour and identity insertions on every
incident physical bond. It then derives a full-support homotopy from
lifted displacement equality, and virtual invariance identifies the actual
state with that of the original path. Native-to-graph conversion preserves
centrality and the incident identities. Central complementary insertions
preserve every local Gram kernel, giving the actual reduced-matrix equality.
Nonvanishing follows from regular injectivity. No detour, homotopy,
factorization, density equality or nonzero-cut assumption is supplied.

Centrality is explicit. Thus all fluxes in abelian groups are covered,
as are central fluxes in nonabelian groups. General noncentral flux pairs
remain under \ghissue{8703}, as does completeness for homotopies constrained
to an arbitrary swept set. The theorem uses full-support homotopy to compare
physical states and does not claim a homotopy avoiding the observation.

Simple connectedness is required of $B(R)$ itself, not inferred from simple
connectedness of $R$ or from endpoint exclusion. An endpoint-free annulus
enclosing one flux is excluded, consistently with the Wilson-loop
obstruction documented in the
\href{https://github.com/LionSR/TNLean/pull/10624}{closed-flux density draft}.
No full-lattice Hamiltonian conclusion under \ghissue{8947} is asserted.
\end{document}
